\section{Capitolo 1: Strutture Dati (Array, Set, Alberi)}

\begin{esercizio}{Es. 1: Il "Flatten" Manuale (Array Annidati)}
\textbf{TRACCIA:} Scrivere una funzione ricorsiva `appiattisci(arr)` che prende un array che può contenere numeri o altri array (annidamento infinito) e restituisce un unico array piatto con tutti i numeri.
\textit{Non usare `flat()` nativo.}

\textbf{Esempio:} `[1, [2, [3, 4]], 5]` $\to$ `[1, 2, 3, 4, 5]`
\end{esercizio}

\begin{jsbox}
function appiattisci(arr) {
    let risultato = [];
    
    for (let i = 0; i < arr.length; i++) {
        let el = arr[i];
        
        if (Array.isArray(el)) {
            // Se è un array, chiamata ricorsiva!
            let sottoArrayPiatto = appiattisci(el);
            // Concateno i risultati
            for (let j = 0; j < sottoArrayPiatto.length; j++) {
                risultato.push(sottoArrayPiatto[j]);
            }
        } else {
            // Se è un numero, lo aggiungo
            risultato.push(el);
        }
    }
    return risultato;
}
\end{jsbox}

\begin{esercizio}{Es. 2: Multi-Insieme da Stringa}
\textbf{TRACCIA:} Data una stringa, creare un oggetto che rappresenti il "Multi-Insieme" dei caratteri (conteggio frequenze), ma ignorando gli spazi.

\textbf{Esempio:} `"a b a"` $\to$ `\{ "a": 2, "b": 1 \}`
\end{esercizio}

\begin{jsbox}
function contaCaratteri(testo) {
    let frequenze = {};
    
    for (let i = 0; i < testo.length; i++) {
        let char = testo[i];
        if (char === " ") continue; // Salta spazi
        
        if (frequenze[char]) {
            frequenze[char]++;
        } else {
            frequenze[char] = 1;
        }
    }
    return frequenze;
}
\end{jsbox}

\begin{esercizio}{Es. 3: Alberi Binari - Somma Foglie}
\textbf{TRACCIA:} Dato un albero binario (oggetti con `val`, `sx`, `dx`), restituire la somma dei valori delle sole \textbf{foglie} (nodi senza figli).

\textbf{Struttura:} `Nodo = { val: number, sx: Nodo | null, dx: Nodo | null }`
\end{esercizio}

\begin{jsbox}
function sommaFoglie(nodo) {
    if (!nodo) return 0;
    
    // È una foglia? (Nessun figlio)
    if (!nodo.sx && !nodo.dx) {
        return nodo.val;
    }
    
    // Altrimenti somma ricorsivamente i sottoalberi
    return sommaFoglie(nodo.sx) + sommaFoglie(nodo.dx);
}
\end{jsbox}

\begin{esercizio}{Es. 4: Intersezione di N Array}
\textbf{TRACCIA:} Scrivere una funzione che accetta un array di array (matrice) e restituisce un array contenente solo gli elementi presenti in \textbf{tutti} i sotto-array.

\textbf{Esempio:} `[[1, 2, 3], [2, 3, 4], [2, 5, 3]]` $\to$ `[2, 3]`
\end{esercizio}

\begin{jsbox}
function intersezioneTutti(matrice) {
    if (matrice.length === 0) return [];
    
    // Parto dai candidati del primo array
    let candidati = matrice[0];
    
    // Per ogni array successivo, filtro i candidati
    for (let i = 1; i < matrice.length; i++) {
        let arrayCorrente = matrice[i];
        
        // Tengo solo chi esiste anche qui
        candidati = candidati.filter(x => arrayCorrente.includes(x));
    }
    
    return candidati;
}
\end{jsbox}
